% intro.tex (v14, 2026-10-20): four paragraphs; the gap and the contributions stated without rhetoric.
The active layer is the near-surface ground above permafrost that thaws each summer and refreezes each winter. Its thickness (ALT) sets the depth of seasonal thaw and controls ground stability, near-surface hydrology and the amount of soil organic carbon exposed to decomposition\cite{Hjort2018,Schuur2015,Biskaborn2019}. ALT responds to the thermal regime of individual years, so that interannual variability is an intrinsic part of the signal recorded by long-term monitoring programmes such as the Circumpolar Active Layer Monitoring (CALM) network\cite{Nelson1997,Streletskiy2025}. Observations, however, are sparse and clustered: most measurements originate from a small number of intensive field campaigns and long-term grids, and many sites are visited only once\cite{Moore2025,Veremeeva2025}.

Several gridded ALT products at about 1 km resolution now cover the northern permafrost region. Process-based estimates such as ESA CCI Permafrost solve the ground heat equation forced by remotely sensed land-surface temperature\cite{Westermann2024,Westermann2025}, whereas statistical and machine-learning (ML) products relate site observations to climate, terrain and soil predictors\cite{Aalto2018,Ran2022a,Liu2024,Wei2026}. Three limitations restrict their use for year-specific applications. First, most products are trained on multi-year mean predictors or on site-years pooled across many years, so that a product value is not tied to the climate of a stated year. Second, products are evaluated with different site sets and validation schemes; reported errors range from a random ten-fold root-mean-square error (RMSE) of 21.6 cm\cite{Liu2024} to a leave-one-site-out RMSE of 60.5 cm\cite{Wei2026} and cannot be compared with each other. Third, random cross-validation of spatially clustered observations overstates the accuracy of ML models at new locations\cite{Roberts2017,Ploton2020,MeyerPebesma2022,Wadoux2021}, which hides how a product performs in a region without observations, that is, where it is most needed.

Physically based models offer a complementary route. The Stefan solution relates thaw depth to the thawing degree-days (TDD) of the season through a single coefficient that depends on soil thermal properties and moisture\cite{Stefan1891,Riseborough2008,Nelson1997}. It reproduces the seasonal progression of thaw and responds to the climate of each year, but one coefficient per region leaves systematic residuals associated with soil organic matter and terrain. Physics-guided ML, in which a physical model carries the dominant response and a learner corrects what the physics omits, has improved accuracy and extrapolation in several environmental applications\cite{Willard2022,Read2019,Jia2021}, and in Alaska the Stefan model retained its accuracy on held-out sites better than a random forest\cite{Gautam2025}. Whether such a combination improves on the Stefan model when the climate input is matched to the observation year, whether it transfers to a region without labels, how many local labels a new region requires, and how it compares with the published products in the same cells and years has not been tested.

Here we address these questions with a site-year design, in which each ALT observation is an observation site in a stated year and each climate input is computed from ERA5-Land for that year and accumulated to the observation day\cite{munozsabater2021_era5_land}. Using \nncyAll{} site-year observations from 2010--2024 in Alaska, the Lena Delta and north-western Canada, we compare the Stefan model, direct ML and a physics-residual workflow in which gradient boosting learns the residual of a recalibrated Stefan model. All methods and five published products are scored against the same observations under the same validation designs, including the hold-out of an entire region and a test in eight external CALM regions, which is what makes their accuracies comparable. The study (i) quantifies the effect of the year-specific climate input, (ii) identifies the conditions under which a learner improves on the physics and the conditions under which the physical anchor preserves transferability, (iii) establishes how many local labels a new region needs, and (iv) provides year-labelled 1 km ALT maps for 2011--2024 with an area of applicability and prediction intervals, together with the evaluation protocol.
